\documentclass[10pt,a4paper]{amsart}
\usepackage[margin=19mm]{geometry}
\usepackage{amsmath,amssymb,mathtools}
\usepackage[scr=rsfs,cal=euler]{mathalfa}
\usepackage{xfrac}
\usepackage[unicode,hidelinks,bookmarks=false]{hyperref}
\newcommand{\ie}{i.e.\ }
\newcommand{\cf}{cf.\ }
\newcommand{\resp}{respectively\ }
\newcommand{\Subsection}{\subsection}
\providecommand{\nobreakdash}{\nobreak\hbox{-}}
\setlength{\emergencystretch}{2em}
\multlinegap0pt
\usepackage{mathtools}
\usepackage[normalem]{ulem}
\newtheorem{proposition}{Proposition}
\usepackage{cleveref}
\crefname{proposition}{Proposition}{Propositions}
\newcommand{\F}{\mathbb{F}}
\newcommand{\dd}{\mathbf{d}}    % the architecture dimensions vector
\newcommand{\qbinom}[2]{\genfrac{[}{]}{0pt}{}{#1}{#2}_q} 
\title{Expressivity of Shallow Neural Networks Over Finite Fields}
\author{Maksym Zubkov, Carol Wu, Shiwei Yang, Param Mody, Yifei Chen}
\date{Source: arXiv:2607.17090v1 (CC BY 4.0)}
\begin{document}
\maketitle

\noindent\emph{Source and licence.} Expressivity of Shallow Neural Networks Over Finite Fields. Version arXiv:2607.17090v1 (CC BY 4.0), distributed by arXiv under the Creative Commons Attribution 4.0 International licence, http://creativecommons.org/licenses/by/4.0/, as recorded in the arXiv licence metadata for this exact version and shown on the abstract page https://arxiv.org/abs/2607.17090v1. Full licence notice: \texttt{licenses/arxiv-2607.17090-CC-BY-4.0.txt}.\par\smallskip

\noindent\textit{Editorial scope.} Counted unit: the article's proposition (n\_2\_k)-r>1 (source0.tex lines 600--606). The article's complete original proof of it is delivered with it (source0.tex lines 607--647, one \texttt{proof} environment of 41 lines). No mathematics is written, repaired or re-derived: the statement and the proof are the article's own lines, in the article's own notation, and no retained line is substituted or re-flowed.  Retained uncounted as the source-local context the proof needs: lines 150--153.  Nothing was omitted from any retained span, so the package carries no omission record.  No retained line contains a citation, so the wrapper carries no bibliography and no bibliography entry.  No retained line contains a figure environment, an image-inclusion command or drawing code, so no image is delivered and the package declares no asset dependency.  Editorial formatting only, in addition: an AMS \texttt{amsart} preamble that restates the article's own declarations and macros used by the retained lines, namely the environments \texttt{proposition} on separate counters, together with the standard packages those lines need.  The retained environments print in this document as Proposition 1 in order of appearance, whereas the article prints the same items with the numbering of its own full document.  The complete native source of the article as distributed in the arXiv e-print for this exact version is bundled unchanged as \texttt{sources/205531/source0.tex}.
    \begin{equation}
    \label{eq: k x n matrices of rank m}
        N_{q}(n,k,m)=\qbinom{k}{m}(q^n-1)(q^n-q)\dots(q^n-q^{m-1}).
    \end{equation}

    \begin{proposition}\label{prop: (n_2_k)-r>1}
        For $\dd = (n,2,k)$, $p \nmid r$, $r > 1$,
        \begin{equation} \label{eq: (n_2_k)-r>1}
            \left|\overline{\mathcal{P}_{\dd, r}}\right| = \left|\overline{\mathcal{P}_{(n,2,1), r}}\right| \left|\mathbb{P}^{k-1}\right| + \binom{\left|\overline{\mathcal{P}_{(n,1,1), r}}\right|}{2}\frac{N_{q}(2,k,2)}{q-1}
        \end{equation}
        where $N_q(n,k,m)$ represents the number of $k\times n$ matrices of rank $m$ over $\F_q$ (see~\Cref{eq: k x n matrices of rank m}).
    \end{proposition} 

    \begin{proof}
        The elements of $\overline{\mathcal{P}_{\dd, r}}$ are projective equivalence classes of ordered $k$-tuples of $n^{\otimes r}$ symmetric tensors $\left[A_1:...:A_k\right]$ of the form 
        $$A_i = b_{i1} L_1^{\otimes r} + b_{i2} L_2^{\otimes r}$$
        where $\left[L_1^{\otimes r}\right], \left[L_2^{\otimes r}\right] \in \overline{\mathcal{P}_{(n,1,1),r}}$ are equivalence classes of rank-1 $n^{\otimes r}$ symmetric tensors. 
    
        We consider $L_1^{\otimes r}$, $L_2^{\otimes r}$ as elements of a vector space over $\mathbb{F}_q$. Then, since we have two basis elements $L_1^{\otimes r}$ and $L_2^{\otimes r}$, $\left[A_1:...:A_k\right]$ spans 1 or 2 dimensions. 
    
        If the tensors in $\left[A_1:...:A_k\right]$ span 1 dimension, then, in fact, all $A_i$ are of the form
        $A_i = b_iA$
        for some $A \in \mathcal{P}_{(n,2,1),r}$. The choice of $\left[A\right]$ and parameterization $\left[B\right] = \left[b_1, b_2, ..., b_k\right] \in \mathbb{P}^{k-1}$ uniquely determines $\left[A_1:...:A_k\right]$. Hence, 
        $$\{\text{1-dimensional} \left[A_1:...:A_k\right]\} \cong \overline{\mathcal{P}_{(n,2,1), r}} \times \mathbb{P}^{k-1}$$ 
        and we obtain the term
        $$\left|\overline{\mathcal{P}_{(n,2,1), r}}\right| \left|\mathbb{P}^{k-1}\right|.$$
    
        In the case of a 2-dimensional span by $\left[A_1:...:A_k\right]$, we begin by arguing the choice of basis pair $\left[L_1^{\otimes r}\right],\left[L_2^{\otimes r}\right]$ partitions the set of 2-dimensional $\left[A_1:...:A_k\right]$ disjointly.
        
        For the sake of contradiction, suppose some 2-dimensional $\left[A_1:...:A_k\right]$ may be constructed by two different basis pairs $\left[L_1^{\otimes r}\right],\left[L_2^{\otimes r}\right]$ and $\left[L_3^{\otimes r}\right],\left[L_4^{\otimes r}\right]$. Since $\left[A_1:...:A_k\right]$ is 2-dimensional, the $A_i$ form an alternative spanning set for the space spanned by $L_1^{\otimes r}$ and $L_2^{\otimes r}$ (or $ L_3^{\otimes r}, L_4^{\otimes r}$) and we can therefore recover the original basis $L_1^{\otimes r},L_2^{\otimes r}$ and $ L_3^{\otimes r}, L_4^{\otimes r}$ as a linear combination of the $A_i$'s. Consequently, $L_3^{\otimes r}, L_4^{\otimes r}$ are both in the span of $L_1^{\otimes r},L_2^{\otimes r}$, and the two pairs form the same plane in $S_r(\F_q^n)$. Note that for the pairs to be distinct, at least one of these linear combinations of $L_1^{\otimes r},L_2^{\otimes r}$ is non-trivial. More precisely, there exists $c_1,c_2 \neq 0$ such that
        $$L_j^{\otimes r} = c_1L_1^{\otimes r} + c_2L_2^{\otimes r}$$
        where $j = 3$ or $4$.

        Because $\left[L_1^{\otimes r}\right] \neq \left[L_2^{\otimes r}\right]$, we have $L_1^{\otimes r}$, $L_2^{\otimes r}$ are linearly independent, and consequently so are $L_1$ and $L_2$. If we apply a change of basis such that $L_1 = (1, 0, 0, ..., 0)$ and $L_2 = (0, 1, 0, ..., 0)$ are standard basis vectors, then $c_1L_1^{\otimes r} + c_2L_2^{\otimes r}$ is an $n^{\otimes r}$ tensor with one $c_1$ and one $c_2$ along the diagonal and 0's everywhere else. Since $r > 1$, therefore this tensor is of symmetric CP-rank 2. In other words, there does not exist an $L_3$ such that $L_3^{\otimes r} = c_1L_1^{\otimes r} + c_2L_2^{\otimes r}$ where $c_1, c_2 \neq 0$. Therefore, we have a contradiction, and in fact, our choice of distinct $\left[L_1^{\otimes r}\right],\left[L_2^{\otimes r}\right]$ will partition the set of 2-dimensional $\left[A_1:...:A_k\right]$'s disjointly.
        
        Fixing a choice of distinct $\left[L_1^{\otimes r}\right]$ and $\left[L_2^{\otimes r}\right] \in \overline{\mathcal{P}_{(n,1,1),r}}$, $\left[A_1:...:A_k\right]$ is 2-dimensional if and only if the coefficients as vectors $\left(\begin{array}{c} b_{i1} \\ b_{i2}\end{array}\right)$ span two dimensions. That is, each possible $\left[A_1:...:A_k\right]$ is uniquely parameterized by a choice of basis pair $\left[L_1^{\otimes r}\right], \left[L_2^{\otimes r}\right] \in \overline{\mathcal{P}_{(n,1,1),r}}$ and some rank-2 matrix
        \[
        \left[(B)\right] = \left[\left(\begin{array}{cccc}
            b_{11} & b_{21} & \hdots & b_{k1} \\
            b_{12} & b_{22} & \hdots & b_{k2} 
        \end{array}\right)^T\right] \in \mathbb{P}(M^{2 \times k}).
        \]
        There are $\binom{\left|\overline{\mathcal{P}_{(n,1,1), r}}\right|}{2}$ ways to choose $\left[L_1^{\otimes r}\right],\left[L_2^{\otimes r}\right]$. Fixing $\left[L_1^{\otimes r}\right],\left[L_2^{\otimes r}\right]$, the number of choices of matrix projective point $\left[(B)\right]$ is $\frac{N_{q}(2,k,2)}{q - 1} = \frac{(q^k - 1)(q^k - q)}{q-1}$. Hence, we obtain our second term, the count of 2-dimensional $\left[A_1:...:A_k\right]$,
        \[
            \binom{\left|\overline{\mathcal{P}_{(n,1,1), r}}\right|}{2} \frac{N_{q}(2,k,2)}{q - 1}.
        \]
        Therefore,
        \begin{align*}
            \left|\overline{\mathcal{P}_{\dd, r}}\right| 
            &= \underbrace{\left|\overline{\mathcal{P}_{(n,2,1),r}}\right| \cdot \left|\mathbb{P}^{k-1}\right|}_{\text{$1$-dimensional span}}
            + \underbrace{\binom{\left|\overline{\mathcal{P}_{(n,1,1),r}}\right|}{2} \cdot \frac{N_q(2,k,2)}{q-1}}_{\text{$2$-dimensional span}}.
        \end{align*}
        as desired.
    \end{proof}

\end{document}
